\documentclass{elsarticle}
\usepackage{amsmath}
\usepackage{amssymb}
\usepackage{amsthm}
\usepackage{hyperref}
\usepackage{tikz}

\title{A second order unstructured finite volume method for incompressible Navier-Stokes equations using face based constrained least-squares procedure}
\author{Dasika Sunder}
\date{}

\begin{document}
	
	\begin{frontmatter}


	\begin{abstract}
		This paper discusses the use of face based constrained least-squares procedure for solving incompressible Navier-Stokes equations. Although reconstruction is performed on faces of the mesh, the underlying finite volume method is still cell centered.  This arrangement of reconstruction is particularly advantageous in implementation     
	\end{abstract}
	
	\end{frontmatter}

	
	\section{Methodology}
	
	
	
	\subsection{Constrained Least Squares}
	The constrained least squares is \cite{boyd2018introduction} used for reconstruction \cite{gradient-calculation-methods} to solve incompressible Navier-Stokes \cite{KIM2000411}
	
	\subsection{Reconstruction for interior faces}
	\noindent The constrained least-squares system for reconstruction on an interior face of the mesh is
	\begin{equation}
		\begin{bmatrix}
			2(N_c-2) & 2\sum \Delta x_i & 2\sum \Delta y_i & 1 & 1\\
			2\sum \Delta x_i & 2\sum \Delta x_i^2 & 2\sum \Delta x_i \Delta y_i & \Delta x_1 & \Delta x_2 \\
			2\sum \Delta y_i & 2\sum \Delta x_i \Delta y_i & 2\sum \Delta y_i^2 & \Delta y_1 & \Delta y_2 \\
			1 & \Delta x_1 & \Delta y_1 & 0 & 0 \\
			1 & \Delta x_2 & \Delta y_2 & 0 & 0 \\
		\end{bmatrix}
		\begin{bmatrix}
			\phi_0 \\
			\phi_x \\
			\phi_y \\
			\hat{z}_1 \\ 
			\hat{z}_2
		\end{bmatrix}	
		= 
		\begin{bmatrix}
			2\sum \phi_i \\
			2\sum \Delta x_i \phi_i  \\ 
			2\sum \Delta y_i \phi_i \\
			\phi_1 \\
			\phi_2 \\ 
		\end{bmatrix}	
	\end{equation}  
	
	\subsection{Reconstruction for faces on Dirichlet boundaries}
	
	\begin{equation}
		\begin{bmatrix}
			2(N_c-2) & 2\sum \Delta x_i & 2\sum \Delta y_i & 1 & 1\\
			2\sum \Delta x_i & 2\sum \Delta x_i^2 & 2\sum \Delta x_i \Delta y_i & \Delta x_1 & 0 \\
			2\sum \Delta y_i & 2\sum \Delta x_i \Delta y_i & 2\sum \Delta y_i^2 & \Delta y_1 & 0 \\
			1 & \Delta x_1 & \Delta y_1 & 0 & 0 \\
			1 & 0 & 0 & 0 & 0 \\
		\end{bmatrix}
		\begin{bmatrix}
			\phi_0 \\
			\phi_x \\
			\phi_y \\
			\hat{z}_1 \\ 
			\hat{z}_2
		\end{bmatrix}	
		= 
		\begin{bmatrix}
			2\sum \phi_i \\
			2\sum \Delta x_i \phi_i  \\ 
			2\sum \Delta y_i \phi_i \\
			\phi_1 \\
			\phi_b \\ 
		\end{bmatrix}	
	\end{equation} 
	
	\subsection{Reconstruction for faces on Neumann boundaries}
	
	\begin{equation}
		\begin{bmatrix}
			2(N_c-2) & 2\sum \Delta x_i & 2\sum \Delta y_i & 1 & 0\\
			2\sum \Delta x_i & 2\sum \Delta x_i^2 & 2\sum \Delta x_i \Delta y_i & \Delta x_1 & n_x \\
			2\sum \Delta y_i & 2\sum \Delta x_i \Delta y_i & 2\sum \Delta y_i^2 & \Delta y_1 & n_y \\
			1 & \Delta x_1 & \Delta y_1 & 0 & 0 \\
			0 & n_x & n_y & 0 & 0 \\
		\end{bmatrix}
		\begin{bmatrix}
			\phi_0 \\
			\phi_x \\
			\phi_y \\
			\hat{z}_1 \\ 
			\hat{z}_2
		\end{bmatrix}	
		= 
		\begin{bmatrix}
			2\sum \phi_i \\
			2\sum \Delta x_i \phi_i  \\ 
			2\sum \Delta y_i \phi_i \\
			\phi_1 \\
			\phi_b \\ 
		\end{bmatrix}	
	\end{equation} 
	
	All summations in the matrix above are from $i = 3$ to $N_c$
	
	\begin{figure}
		\centering
		
		\begin{tikzpicture}
			%\draw[help lines] grid(3,3);
			\draw (0,0) -- (4,0) -- (2,{2*sqrt(3)}) -- cycle;
		\end{tikzpicture}
		
		\caption{Arrangement of variables in a cell}
		
	\end{figure}
	

	
	
	
	\bibliographystyle{plain}
	\bibliography{references}

	
	
\end{document}